\documentclass[11pt]{article}

\usepackage[margin=1in]{geometry}
\usepackage{amsmath, amssymb}
\usepackage{graphicx}
\usepackage{float}
\usepackage[hidelinks]{hyperref}
% Long file:line paths in \texttt would otherwise overflow the margin.
\usepackage[htt]{hyphenat}
\sloppy

\newcommand{\dt}{\Delta t}
\newcommand{\uh}{u_h}
\newcommand{\Ecvx}{E_{\mathrm{cvx}}}
\newcommand{\Eccv}{E_{\mathrm{ccv}}}

\title{Allen--Cahn: background, discretisation, and experiments}
\date{\today}

\begin{document}
\maketitle

\noindent
A self-contained account of the Allen--Cahn part of the project: the energy, the gradient flow,
the discretisation as it is actually implemented, the two time-stepping schemes and the two
physics losses they induce, the autodiff-through-time question, and a summary of every experiment
we have run. Everything is checked against the source; pointers are \texttt{path:line} at the time
of writing. The companion document \texttt{notes/ac\_autoregressive/ac\_autoregressive.tex} carries
the longer derivations (numerical regime analysis, conditioning, mesh-consistent scaling); this
document is the summary and the experimental record.

\section{The energy}

Allen--Cahn is the $L^2$ gradient flow of the Ginzburg--Landau free energy. On $\Omega=(0,1)^2$
with homogeneous Dirichlet conditions,
\begin{equation}\label{eq:energy}
   E(u) \;=\; \frac{a^2}{2}\int_\Omega |\nabla u|^2\,\mathrm dx
            \;+\; \epsilon^2\!\int_\Omega W(u)\,\mathrm dx,
   \qquad
   W(s) = \tfrac14\,(s^2-1)^2 .
\end{equation}
The first term penalises interfaces, the double well $W$ has minima at $s=\pm1$ and forces the
field towards the two pure phases. The competition sets an interface width $\ell \sim a/\epsilon$.

Two conventions that are easy to get wrong and are fixed throughout the code and this project:
\begin{itemize}
   \item $\epsilon$ multiplies the \emph{reaction}, not the Laplacian, and it enters
         \emph{squared}. Large $\epsilon$ means a stiff, sharply-separating problem. This follows
         TensorMesh's \texttt{examples/diffusion/allen-cahn} and is the opposite of the more common
         $\epsilon\to0$ sharp-interface convention.
   \item The diffusion coefficient is $a^2$; we use $a=1$ in every experiment, so $\epsilon$ alone
         controls stiffness.
\end{itemize}

\section{The gradient flow}

The $L^2$ gradient flow $\partial_t u = -\,\delta E/\delta u$ of \eqref{eq:energy} is
\begin{equation}\label{eq:ac-pde}
   \partial_t u \;=\; a^2\Delta u \;+\; \epsilon^2\,u(1-u^2),
   \qquad u|_{\partial\Omega}=0,
\end{equation}
since $-\delta E/\delta u = a^2\Delta u - \epsilon^2 W'(u)$ and $W'(s)=s^3-s$. The structural
consequence we lean on repeatedly is that $E$ is a Lyapunov functional,
\begin{equation}\label{eq:dissipation}
   \frac{\mathrm d}{\mathrm dt}E(u(t)) \;=\; -\,\|\partial_t u(t)\|_{L^2}^2 \;\le\; 0 ,
\end{equation}
so \textbf{the energy must be non-increasing along any correct trajectory}. We use this as a
diagnostic for learned rollouts: an energy curve that climbs is a blow-up, regardless of what the
error metric says (\texttt{ACProblem.energy}, \texttt{tensorpils/physics.py:282}).

\section{Discretisation: nodal interpolation of the nonlinearity}

Space is discretised with bilinear $Q_1$ elements on a structured quadrilateral grid of the unit
square (\texttt{structured\_quad\_mesh}, \texttt{tensorpils/meshing.py:24}); the stiffness $A$ and
mass $M$ are assembled once by TensorMesh with Gauss quadrature of order 4 per axis
(\texttt{physics.py:64--65}). This is the same $A$, $M$, and Dirichlet masking as the Poisson
setting: the boundary condition is imposed by projection, not condensation.

The nonlinearity is treated by \textbf{nodal (group / product) interpolation}: rather than
integrating $W'(\uh)$ exactly, we interpolate the nonlinear function of the nodal values into the
same $Q_1$ space,
\begin{equation}\label{eq:group-fem}
   \int_\Omega W'(\uh)\,\phi_i\,\mathrm dx
   \;\approx\;
   \int_\Omega \Big(\textstyle\sum_j W'(u_j)\,\phi_j\Big)\phi_i\,\mathrm dx
   \;=\;\big[M\,W'(\mathbf u)\big]_i ,
\end{equation}
with $W'(\mathbf u)$ applied entrywise. So the cubic reaction becomes $M(\mathbf u-\mathbf u^3)$
and never requires quadrature of a cubic integrand. The same convention gives the discrete energy
\begin{equation}\label{eq:discrete-energy}
   E_h(\mathbf u) \;=\; \tfrac12 a^2\,\mathbf u^\top A\,\mathbf u
   \;+\; \epsilon^2\,\mathbf 1^\top M\,W(\mathbf u),
\end{equation}
i.e.\ the double well is evaluated nodally and then integrated exactly against the $Q_1$ basis
(\texttt{physics.py:282--298}). Nodal interpolation is what makes the reaction Jacobian a
\emph{diagonal scaling of $M$}, which is exactly what the Newton solve and the preconditioner
below exploit.

\section{Time stepping I: minimizing movement}

\paragraph{Continuous (JKO / De Giorgi).}
The gradient-flow structure gives a variational time discretisation: given $u^n$, define
\begin{equation}\label{eq:mm-continuous}
   u^{n+1} \;=\; \arg\min_{u}\;\Big\{\, E(u) \;+\; \frac{1}{2\dt}\,\|u-u^n\|_{L^2}^2 \,\Big\}.
\end{equation}
The optimality condition is $\delta E/\delta u + (u-u^n)/\dt = 0$, i.e.\ exactly backward Euler
applied to \eqref{eq:ac-pde}. Choosing $u^{n+1}$ as the minimiser rather than as the root of a
residual guarantees $E(u^{n+1}) + \tfrac{1}{2\dt}\|u^{n+1}-u^n\|^2 \le E(u^n)$ and hence
unconditional energy decay. The catch is that $E$ is \emph{not} convex --- the concave part
$-\tfrac{\epsilon^2}{2}\int u^2$ of the double well spoils it --- so \eqref{eq:mm-continuous} is
not a convex problem, and for $\dt$ large relative to $1/\epsilon^2$ the objective can lose its
unique minimiser. This is what the convex--concave split of \S\ref{sec:cc} repairs.

\paragraph{Discrete (what we implement).}
We use the \emph{convex--concave} minimizing-movement objective throughout, so \S\ref{sec:cc}
supplies the split and this section supplies the variational form. Splitting
$E = \Ecvx + \Eccv$ with
\begin{equation}\label{eq:split}
   \Ecvx(u) = \frac{a^2}{2}\!\int|\nabla u|^2 + \frac{\epsilon^2}{4}\!\int u^4
              + \frac{\epsilon^2}{4}|\Omega|,
   \qquad
   \Eccv(u) = -\frac{\epsilon^2}{2}\!\int u^2 ,
\end{equation}
(so that $\Ecvx+\Eccv = E$, using $W(s)=\tfrac14(s^4-2s^2+1)$), and linearising only the concave
part about $u^n$, the step is the unique minimiser of the \emph{convex} objective
\begin{equation}\label{eq:mm-objective}
   J(u) \;=\; \Ecvx(u) \;+\; \big\langle D\Eccv(u^n),\,u\big\rangle
              \;+\; \frac{1}{2\dt}\|u-u^n\|_{L^2}^2,
   \qquad D\Eccv(u^n) = -\epsilon^2 u^n .
\end{equation}
In FEM variables, with the quartic evaluated nodally against the \emph{lumped} mass
$M\mathbf 1$ (\texttt{ACProblem.mm\_objective}, \texttt{physics.py:240--280}),
\begin{equation}\label{eq:mm-discrete}
   J(\mathbf u) \;=\;
     \frac{a^2}{2}\mathbf u^\top A\mathbf u
   \;+\; \frac{\epsilon^2}{4}\Big[(\mathbf u^4)^\top M\mathbf 1 + |\Omega|\Big]
   \;-\; \epsilon^2 (\mathbf u^n)^\top M \mathbf u
   \;+\; \frac{1}{2\dt}(\mathbf u-\mathbf u^n)^\top M(\mathbf u-\mathbf u^n),
\end{equation}
whose gradient is
$\nabla J = M\frac{\mathbf u-\mathbf u^n}{\dt} + a^2A\mathbf u
+ \epsilon^2\operatorname{diag}(M\mathbf 1)\mathbf u^3 - \epsilon^2 M\mathbf u^n$.
Because $J$ is convex with Hessian
$\nabla^2 J = M/\dt + a^2A + 3\epsilon^2\operatorname{diag}(M\mathbf 1)\operatorname{diag}(\mathbf u^2)$,
which is SPD for every $\mathbf u$ and every $\dt$, this is the \textbf{Deep-Ritz analogue for
time stepping}: the loss is an energy to be minimised, not a residual to be squared, and it
inherits $\kappa(\nabla^2 J)$ rather than $\kappa^2$. That is the entire reason it behaves
differently from the least-squares loss.

\paragraph{Small inconsistency worth knowing.}
$J$ uses the \emph{lumped} mass on the quartic ($(\mathbf u^4)^\top M\mathbf 1$) while the
diagnostic energy \eqref{eq:discrete-energy} uses the \emph{consistent} mass
($\mathbf 1^\top M\,W(\mathbf u)$). Both are legitimate nodal-interpolation choices and they agree
to quadrature order, but they are not literally the same functional, so $J$ and $E_h$ should not
be expected to match to machine precision.

\section{Time stepping II: convex--concave (Eyre) splitting}
\label{sec:cc}

\paragraph{Continuous.}
The Eyre split treats the convex part implicitly and the concave part explicitly:
\begin{equation}\label{eq:cc-continuous}
   \frac{u^{n+1}-u^n}{\dt}
   \;=\; a^2\Delta u^{n+1} \;-\; \epsilon^2 (u^{n+1})^3 \;+\; \epsilon^2 u^n .
\end{equation}
Compare fully implicit backward Euler, which puts the whole reaction
$\epsilon^2\big(u^{n+1}-(u^{n+1})^3\big)$ at the new level. The two differ \emph{only in the
linear reaction term}: implicit at $u^{n+1}$ (backward Euler) versus explicit at $u^n$
(convex--concave). The split is unconditionally energy-stable and is the standard scheme for
Allen--Cahn at stiff $\epsilon$; it is exactly the optimality condition of \eqref{eq:mm-objective},
so \emph{the convex--concave step and the minimizing-movement step are the same step}, once via a
residual and once via an objective.

\paragraph{Discrete.}
With nodal interpolation the two residuals are (\texttt{ACProblem.residual},
\texttt{physics.py:205--237}), both scaled by $1/\dt$:
\begin{align}
   R_{\mathrm{BE}}(\mathbf u^n,\mathbf u^{n+1})
   &= M\frac{\mathbf u^{n+1}-\mathbf u^{n}}{\dt} + a^2A\,\mathbf u^{n+1}
      - \epsilon^2 M\big(\mathbf u^{n+1}-(\mathbf u^{n+1})^3\big), \label{eq:res-be}\\[2pt]
   R_{\mathrm{CC}}(\mathbf u^n,\mathbf u^{n+1})
   &= M\frac{\mathbf u^{n+1}-\mathbf u^{n}}{\dt} + a^2A\,\mathbf u^{n+1}
      + \epsilon^2 M\big((\mathbf u^{n+1})^3-\mathbf u^{n}\big). \label{eq:res-cc}
\end{align}
The $1/\dt$ scaling is deliberate: it makes a least-squares loss built on either residual have a
$\dt$-independent gradient scale. The corresponding Newton Jacobians are
\begin{equation}\label{eq:jacobians}
   J_{\mathrm{CC}} = \frac{M}{\dt} + a^2A + 3\epsilon^2 M\operatorname{diag}(\mathbf u^2),
   \qquad
   J_{\mathrm{BE}} = \frac{M}{\dt} + a^2A + \epsilon^2 M\operatorname{diag}(3\mathbf u^2-1).
\end{equation}
$J_{\mathrm{CC}}$ is SPD for every $\mathbf u$ and $\dt$ (the reaction term is a non-negative
diagonal scaling of $M$); $J_{\mathrm{BE}}$ is not, which is why the reference solver defaults to
convex--concave.

\paragraph{The reference trajectory.}
Allen--Cahn has no analytical solution, so the ground truth is a FEM solve: each step drives
$R=0$ by Newton on the interior degrees of freedom, with a \textbf{batched sparse} linear solve
and \textbf{always in \texttt{float64}} regardless of the caller's dtype --- \texttt{fp32} cannot
reach \texttt{newton\_tol}$=10^{-8}$ and leaves Newton spinning to \texttt{newton\_max} every step
(\texttt{ACProblem.fem\_reference}, \texttt{physics.py:380--458}). The Jacobian is reassembled per
Newton step as a value update on the shared sparsity pattern of $L=M/\dt+a^2A$, since the reaction
contributes only a column scaling of $M$ --- again a consequence of nodal interpolation.

Two things follow, and both matter for reading the results:
\begin{itemize}
   \item The learned stepper is measured against the \emph{discrete} convex--concave trajectory,
         not against physical Allen--Cahn. This is well-posed at every $\epsilon$ (the CC scheme is
         unconditionally stable), but at large $\epsilon$ the reference is under-resolved and is no
         longer physical Allen--Cahn.
   \item Data generation and the physics loss use the same integrator by default, so the reference
         \emph{exactly} zeroes the residual the loss penalises (\texttt{ACTrainer},
         \texttt{trainer.py:812}). Decoupling them (\texttt{--ac\_loss\_integrator}) is an
         intentional experiment, not an accident.
\end{itemize}

\section{The two physics losses}

Both losses act on an autoregressive rollout $\mathbf u^0,\dots,\mathbf u^{L-1}$ where
$\mathbf u^0$ is the ground-truth seed frame and the rest are model predictions, with a discount
$\gamma^k$ over rollout steps and a normalisation by $\sum_k\gamma^k$:
\begin{align}
   \mathcal L_{\mathrm{LS}}
     &= \frac{1}{\sum_k \gamma^k}\sum_{k} \gamma^k\;
        \frac12\big\|\,P\,R(\mathbf u^{k},\mathbf u^{k+1})\,\big\|_2^2
        &&\text{(\texttt{ACGalerkinLoss}, \texttt{losses.py:255})} \label{eq:loss-ls}\\
   \mathcal L_{\mathrm{MM}}
     &= \frac{1}{\sum_k \gamma^k}\sum_{k} \gamma^k\;
        J(\mathbf u^{k};\mathbf u^{k+1})
        &&\text{(\texttt{ACMinMovementLoss}, \texttt{losses.py:310})} \label{eq:loss-mm}
\end{align}
with $P=I$ unless preconditioning is requested. The squared residual is \textbf{summed over nodes
and averaged over the batch} (keeping the gradient $O(1)$ rather than $O(1/N)$); $J$ is already a
spatial functional, so no per-node reduction is applied. Note that $\mathcal L_{\mathrm{MM}}$ is an
energy and is routinely \emph{negative} --- in the reference run below it converges to
$\approx-120$ --- so its numerical value is not comparable across losses and only the validation
error is meaningful for model selection.

\paragraph{Preconditioning the least-squares loss.}
$\mathcal L_{\mathrm{LS}}$ has Gauss--Newton Hessian $\approx J^\top J$, so it inherits
$\kappa(J)^2$ --- the same squaring that kills the Poisson least-squares loss. We precondition with
a \textbf{frozen-coefficient} approximation of the Jacobian: setting $\mathbf u^2\equiv 1$ (exact in
the bulk of either phase, wrong only inside the interface) in \eqref{eq:jacobians} gives the
constant, SPD, $u$-independent operator
\begin{equation}\label{eq:frozen}
   J_0 \;=\; a^2 A + c\,M, \qquad c = \frac{1}{\dt} + 3\epsilon^2 ,
\end{equation}
a screened-Poisson / reaction--diffusion operator. $P\approx J_0^{-1}$ is realised by the
\emph{same} geometric-multigrid V-cycle as in the Poisson experiments, simply built on
$a^2A+cM$ instead of $A$ (\texttt{GeometricMultigrid} with \texttt{mg\_a2}/\texttt{mg\_c},
\texttt{preconditioners/multigrid.py:111--125}; wired at \texttt{trainer.py:825--838}). The mass
shift only makes the operator more diagonally dominant, so weighted-Jacobi smoothing is at least as
effective as in the pure Poisson case. Preconditioning applies to the least-squares form only ---
it is meaningless for $J$, whose Hessian is already the un-squared $\nabla^2 J$.

\section{Autodiff through time}

The rollout has \emph{two independent detach knobs}, which is the part that is easy to conflate
(\texttt{RolloutTrainer.\_rollout}, \texttt{trainer.py:622--640}; resolution at
\texttt{trainer.py:820}):
\begin{enumerate}
   \item \textbf{Rollout-input detach.} Whether step $k{+}1$ receives $\mathbf u^{k}$ or
         $\mathbf u^{k}.\mathrm{detach}()$ as \emph{network input}. Detaching makes every stored
         frame a one-step function of its input --- no backpropagation through time at all. This is
         the \emph{pushforward trick} (Brandstetter, Worrall \& Welling, 2022).
   \item \textbf{Loss-coupling detach.} Whether the \emph{previous frame} is differentiated
         \emph{inside} the loss --- the first argument of $R(\mathbf u^k,\mathbf u^{k+1})$, or the
         proximal centre $\mathbf u^k$ in $J$.
\end{enumerate}
The three exposed modes are
\begin{center}
\begin{tabular}{lccl}
\hline
\texttt{--bptt\_mode} & input detach & coupling detach & gradient \\
\hline
\texttt{full\_bptt}   & no  & no  & backprop through the entire rollout \\
\texttt{detach\_prev} & no  & yes & full BPTT, previous frame frozen in the loss \\
\texttt{pushforward}  & yes & yes & one-step \\
\hline
\end{tabular}
\end{center}
The data loss has no coupling term, so \texttt{detach\_prev}$\equiv$\texttt{full\_bptt} there.

\paragraph{Why the coupling detach is not cosmetic.}
For the least-squares residual, $\tfrac12\|R\|^2$ is \emph{self-balancing}: its gradient vanishes in
\emph{both} arguments at $R=0$, so leaving the coupling attached is harmless, and the class default
is not to detach. For minimizing movement it is different: $\nabla_{\mathbf u^k}J$ does \emph{not}
vanish at the correct step --- $J$ is minimised over its second argument only, and the proximal
centre is data. Backpropagating through it therefore rewards moving the \emph{centre}, distorting
the trajectory. The MM class default is to detach. \S\ref{sec:exp-bptt} shows this is worth roughly
two orders of magnitude, and it is the single clearest experimental result we have about autodiff
through time.

\section{Experiments}

All runs: FNO ($16\times16$ modes, width $64$, $5$ layers), Adam with cosine
$10^{-3}\!\to\!10^{-4}$, batch $32$, $n_{\text{train}}/n_{\text{val}}/n_{\text{test}} =
1024/128/256$, \texttt{seed=42}, $a=1$, zero Dirichlet, convex--concave reference. Initial
conditions are multi-frequency,
$u_0 = \frac{\pi}{K^2}\sum_{i,j=1}^{K} a_{ij}(i^2+j^2)^{-1/2}\sin(\pi i x)\sin(\pi j y)$ with
$a_{ij}\sim\mathrm{Unif}[-1,1]$ (amplitude $\approx0.15$), so $K$ controls how rich the initial
data is. Metric is the FEM (mass-matrix) relative $L^2$ against the reference; ``space-time''
averages over the rollout.

\subsection{\texttt{ac\_loss\_comparison} --- does the loss matter at all?}
$64^2$, $\dt=0.0025$, $\epsilon=2$, $K=4$, 10-step rollout, 1000 epochs. Four arms: data-driven,
backward-Euler least squares, convex--concave least squares, minimizing movement. Arm 2 trains the
\emph{BE} residual against \emph{CC} data on purpose, to test whether the physics-loss integrator
matters. \textbf{Outcome:} at $\epsilon=2$ the dynamics are diffusion-dominated (the double well
barely engages) and all four losses land close together --- an informative null. This motivated the
stress test.

\subsection{\texttt{ac\_stresstest} --- sweep the reaction strength}
Same setup, $\epsilon\in\{2,4,8,12,16,20,24,32\}$ (24 runs), $500$ epochs, three losses. Best
space-time relative $L^2$:
\begin{center}
\begin{tabular}{rccc}
\hline
$\epsilon$ & data-driven & least squares & min.\ movement \\
\hline
 2 & 0.024 & 0.022 & 0.013 \\
 4 & 0.025 & 0.004 & 0.005 \\
 8 & 0.011 & 0.004 & 0.004 \\
12 & 0.011 & 0.004 & 0.011 \\
16 & 0.999 & 0.975 & 0.021 \\
20 & 1.000 & 0.991 & 0.056 \\
24 & 1.000 & 0.995 & 0.095 \\
32 & 0.981 & 0.998 & 1.011 \\
\hline
\end{tabular}
\end{center}
Read at face value this is the headline Allen--Cahn story: past $\epsilon\approx12$ both
data-driven and bare least squares collapse to $\approx100\%$ error while minimizing movement
degrades gracefully until it too fails at $\epsilon=32$. \textbf{But see
\S\ref{sec:exp-inconsistency} before using these numbers.}

\subsection{\texttt{compare\_bptt} --- autodiff modes $\times$ reaction strength}
\label{sec:exp-bptt}
$64^2$, $\dt=0.0025$, $K=4$, $500$ epochs, $\epsilon\in\{4,\dots,32\}$, three losses $\times$ their
modes $=64$ runs. Best space-time relative $L^2$:
\begin{center}
\begin{tabular}{llcccccccc}
\hline
loss & mode & 4 & 8 & 12 & 16 & 20 & 24 & 28 & 32 \\
\hline
data & \texttt{full}   & 0.010 & 0.011 & 0.007 & 0.009 & 0.010 & 0.012 & 0.061 & 0.038 \\
data & \texttt{push}   & 0.009 & 0.007 & 0.006 & 0.005 & 0.006 & 0.009 & 0.010 & 0.017 \\
LS   & \texttt{full}   & 0.023 & 0.004 & 0.004 & 0.006 & 0.012 & 0.995 & 0.997 & 0.998 \\
LS   & \texttt{detach} & 0.004 & 0.004 & 0.006 & 0.008 & 0.015 & 0.028 & 0.044 & 1.000 \\
LS   & \texttt{push}   & 0.004 & 0.004 & 0.005 & 0.012 & 0.015 & 0.039 & 0.034 & 0.035 \\
MM   & \texttt{full}   & 0.408 & 0.964 & 1.018 & 2.854 & 2.394 & 2.091 & 1.736 & 1.515 \\
MM   & \texttt{detach} & 0.005 & 0.004 & 0.004 & 0.007 & 0.007 & 0.009 & 0.019 & 1.000 \\
MM   & \texttt{push}   & 0.003 & 0.003 & 0.006 & 0.004 & 0.008 & 0.007 & 0.007 & 0.012 \\
\hline
\end{tabular}
\end{center}
Three clean conclusions:
\begin{itemize}
   \item \textbf{Pushforward is best or tied-best for every loss at every $\epsilon$}, and it is
         also the cheapest (no rollout graph, so activation memory is a single forward). This is
         why it is the default in all later experiments.
   \item \textbf{Minimizing movement with full BPTT is catastrophic} ($0.41\to2.85$, i.e.\ worse
         than predicting nothing), while the same loss with the proximal centre detached is the
         best arm in the table. This is the quantitative version of \S6.
   \item Detaching rescues least squares at large $\epsilon$: \texttt{full} collapses from
         $\epsilon=24$, \texttt{detach} survives to $28$, \texttt{push} holds throughout.
\end{itemize}

\subsection{\texttt{long\_rollout} --- extrapolation past the training horizon}
Rolls the best \texttt{compare\_bptt} checkpoints $20$ steps (trained on $10$) and tracks both the
relative $L^2$ and the Ginzburg--Landau energy \eqref{eq:discrete-energy} against the CC reference,
with divergence detection (non-finite frames are marked and the curve truncated). The energy check
is the point: by \eqref{eq:dissipation} a rising $E$ is a blow-up even if the error looks
acceptable. Two-stage by design --- the compute step runs on Euler where the checkpoints live and
emits a small \texttt{metrics.json}; only that and the figures are transferred.

\subsection{\texttt{realistic\_ac} --- the stiff, resolved regime with preconditioning}
\label{sec:exp-realistic}
The current headline experiment: $128^2$, $\epsilon=32$, $\dt=0.01$, $K=4$, 10-step pushforward
rollout, $500$ epochs, trained then rolled out to \textbf{100 steps}. Four arms: data-driven,
minimizing movement, bare least squares, and \textbf{preconditioned least squares} with the
multigrid V-cycle on $a^2A+cM$. Relative $L^2$ against the CC reference on the test set:
\begin{center}
\begin{tabular}{lcccc}
\hline
loss & best space-time & step 10 & step 50 & step 100 \\
\hline
data-driven                & 0.019 & 0.019 & 0.041 & 0.041 \\
minimizing movement        & 0.008 & 0.008 & 0.014 & 0.013 \\
preconditioned LS          & 0.047 & 0.034 & 0.039 & 0.037 \\
bare LS                    & \textbf{1.000} & \multicolumn{3}{c}{---\ (collapsed; excluded from the rollout figure)} \\
\hline
\end{tabular}
\end{center}
This is the cleanest statement of the Allen--Cahn result:
\begin{itemize}
   \item \textbf{Bare least squares fails completely} --- $\approx100\%$ relative error. It does
         \emph{not} produce NaNs; it stays finite and collapses towards the trivial prediction
         (final training loss $0.069$, versus $3.8\times10^{-3}$ preconditioned).
   \item \textbf{Preconditioning rescues it} to within a factor of two of data-driven training,
         using the same V-cycle machinery as Poisson with only the operator changed.
   \item \textbf{Minimizing movement is the best arm overall}, beating even data-driven training,
         consistent with it inheriting $\kappa$ rather than $\kappa^2$.
   \item \textbf{Rollouts stay stable to 100 steps} despite a 10-step training horizon; the error
         plateaus rather than growing, for all three working losses.
\end{itemize}

\subsection{\texttt{realistic\_ac\_K8} / \texttt{realistic\_ac\_K16} --- harder initial data}
Identical protocol, with the multi-frequency initial condition at $K=8$ and $K=16$ instead of
$K=4$; separate output trees, analysis scripts reused unchanged. $K=8$ is the interesting
intermediate difficulty; at $K=16$ \textbf{none of the arms converged}, which bounds how rich the
initial data can be before the whole approach stops working at this resolution and budget. Locally
only three of the four $K=8$ runs are present and $K=16$ results have not been pulled --- worth
completing before any of this goes in the paper.

\subsection{\texttt{bare\_ls\_seeds} --- is the bare-LS failure deterministic?}
Three seeds of the bare (unpreconditioned) least-squares arm with the NaN guard removed, each
writing to its own output directory, to distinguish ``collapses to $\approx100\%$ but stays
finite'' from ``diverges to NaN''. This is the diagnostic that makes the failure claim in
\S\ref{sec:exp-realistic} safe to state.

\section{Open issues and things worth checking}

\paragraph{An unresolved inconsistency between two sweeps.}
\label{sec:exp-inconsistency}
\texttt{ac\_stresstest} and \texttt{compare\_bptt} have \emph{byte-identical} CLI configurations
apart from \texttt{--bptt\_mode}, which \texttt{ac\_stresstest} leaves unset --- and unset resolves
to \texttt{full\_bptt} for the data and Galerkin losses. So \texttt{ac\_stresstest}'s ``data'' and
``LS'' columns should reproduce \texttt{compare\_bptt}'s \texttt{full} rows. They do not:
\begin{center}
\begin{tabular}{lccc}
\hline
$\epsilon=16$, data-driven & final train loss & best space-time rel.\ $L^2$ & run date \\
\hline
\texttt{ac\_stresstest}  & $1.5\times10^{-1}$ & 0.999 & 10 Jul \\
\texttt{compare\_bptt}   & $1.0\times10^{-5}$ & 0.009 & 15 Jul \\
\hline
\end{tabular}
\end{center}
A \emph{supervised} run failing to fit is not a plausible scientific finding, so the
\texttt{ac\_stresstest} large-$\epsilon$ collapse is most likely a stale-code artefact: the two
sweeps are five days apart, their result JSONs differ five-fold in size (schema change), and the
proximal-centre fix \texttt{ee1fdd0} landed the same day as the stress test. \textbf{Treat
\texttt{compare\_bptt} as authoritative and re-run \texttt{ac\_stresstest} before using its numbers
anywhere.} Note that this undercuts the ``data-driven falls apart as $\epsilon$ grows'' claim in
the \texttt{ac\_stresstest} README --- \texttt{compare\_bptt} shows data-driven holding to
$\epsilon=32$.

\paragraph{Stale READMEs.} \texttt{experiments/allen\_cahn/realistic\_ac/README.md} describes
$256^2$ and $\epsilon=64$ throughout, but \texttt{sweep.sbatch} actually runs $128^2$ and
$\epsilon=32$ (the numbers reported in \S\ref{sec:exp-realistic}). The
\texttt{ac\_stresstest} README lists six $\epsilon$ values; eight were run.

\paragraph{The preconditioned-LS gap.} Preconditioned LS ($0.047$) is still $6\times$ worse than
minimizing movement ($0.008$) and $2.5\times$ worse than data-driven. The frozen-coefficient
approximation $\mathbf u^2\equiv1$ is exact in the bulk and wrong exactly on the interface, which
is where all the dynamics are --- so this gap is expected, but it has not been measured against a
non-frozen (per-step, per-sample) preconditioner. That is the obvious next experiment, and it is
the one that would let us claim preconditioning closes the gap for Allen--Cahn as cleanly as it
does for Poisson.

\paragraph{Not yet done.} No mesh-refinement study for Allen--Cahn (the $h^{-4}$-versus-$h^{-2}$
story is Poisson-only so far); no infinite-data-limit run (streaming datasets exist for Poisson
only); no seed replication for anything except bare LS; and the $\epsilon$-sweep has never been run
with the preconditioned loss included, so we cannot yet say at which $\epsilon$ preconditioning
stops rescuing least squares.

\end{document}
